\section{Discussion}
\label{sec:discussion}

\subsection{Why the belief exists}
The physics we model does not make silent passages generally worse. At equal mass the silent advantage at a
neighbour is a median factor of 1.35, often below the just-noticeable ratio. As emitted, a silent passage is the
stronger one in 47\% of pairs in the analytic model and in 21\% (0--72\%) of bystander pairs in the simulated scenes,
where on a sofa it misses both noses on its first pass.

The belief is nevertheless documented across languages. Theorem~3 accounts for it without any physical difference: a
passage enters memory if it is heard or smelled, so every remembered silent passage was smelled, whereas audible
passages are remembered whether or not they smelled. With literature observer models this selection turns a truth of
0.32--0.51 into a remembered belief of 0.67--0.92 (Table~\ref{tab:observers}). Three documented memory biases act in
the same direction: distinctive pairings are over-remembered \cite{hamilton1976}, vivid episodes are easy to recall
\cite{tversky1973}, and judgements built from biased samples are rarely corrected \cite{fiedler2000,denrell2005}. The
phrase itself can also pass the belief on without any experience \cite{chapman1967}; documentary evidence cannot
separate the two routes (Figure~\ref{SI-fig:culture}).

\subsection{The releaser}
The releaser feels every release, so their noticing does not depend on smell and Theorem~3 does not apply to them.
Three findings bear on how they judge their own passages.
\begin{itemize}
  \item People rate their own flatus as less disgusting than a stranger's: on a 1--5 scale the item ``silent but
  smells strongly'' is rated 2.51 for one's own flatus and 3.71 for a stranger's \cite{liuzza2016,stevenson2005}.
  This lowers the unpleasantness of the releaser's own passages, not the contrast between silent and audible ones.
  \item The self-advantage appears even when people do not recognise their own odour \cite{mcburney1976}.
  \item Believing an odour is one's own shifts its hedonic rating \cite{hold1977}, so expectation is the one route by
  which a releaser could share the belief.
\end{itemize}

\subsection{Where the odour comes from}
\label{sec:disc-chem}
Two facts the transport model takes as given have chemical explanations in the literature; we did not test them.

\emph{Why sulfur concentration does not follow volume.}
\begin{itemize}
  \item \ce{H2S} comes from sulfate-reducing bacteria and from cysteine degradation, whose genes are widespread in the
  gut microbiome \cite{gibson1993,braccia2021}. Sulfate reducers outcompete methanogens for \ce{H2} by affinity
  ($K_s \approx 1$ against 6\,\textmu mol\,L$^{-1}$), and sulfide and methane production tend to exclude each other
  \cite{gibson1993}.
  \item The caecal mucosa of rats removes both gases quickly: methanethiol is demethylated to \ce{H2S}, and \ce{H2S} is
  oxidised mainly to thiosulfate \cite{levitt1999}. Rhodanese is the main detoxifying enzyme in human colon
  \cite{picton2002}, and SELENBP1 oxidises methanethiol with an apparent $K_m$ of 4.8\,nM \cite{pol2018}.
  \item Only free sulfide can enter the gas, and it is less than 5\% of faecal total sulfide \cite{tangerman2009}.
\end{itemize}
The odorants are thus trace products held at a balance of production and fast consumption, whereas passage volume is
set by \ce{N2}, \ce{CO2}, \ce{H2} and \ce{CH4}. We infer that this is why sulfur concentration and passage volume are
uncorrelated \cite{suarez1997}. Diet moves the balance: a meat-rich diet raised faecal sulfide from 0.22 to
3.38\,mmol\,kg$^{-1}$ in five volunteers \cite{magee2000}, and the response to animal and plant diets differs between
people \cite{teigen2022}.

\emph{Why methanethiol smells foul.} The human receptor OR2T11 responds to short thiols, including methanethiol, in a
copper-dependent way, but not to \ce{H2S} \cite{li2016}. Human nasal mucus contains about 2.5\,\textmu M copper
\cite{block2017}, below the 30\,\textmu M added in the receptor assays, so the receptor's role in vivo is plausible
but not shown. Other receptors respond to sulfur ligands without copper \cite{duan2012,haag2019,zhang2018}. A
receptor that responds to methanethiol and not to \ce{H2S} would make methanethiol the characteristic note of flatus
independently of thresholds; we infer this, it has not been tested.

\subsection{Predictions}
Theorem~3 predicts that releasers hold the belief less firmly than bystanders and that people without a sense of
smell never remember a silent passage. Because masking removes the quieter audible passages, the effect of background
noise on the belief should be ambiguous.

\subsection{When a silent release is worse}
At equal mass, at the releaser's own nose, in still air and within a few seconds, a silent release is perceptibly
stronger than a loud one (Figure~\ref{fig:uq}a). In the elevator, a silent release is also stronger for the person
standing behind a front-row releaser, who smells it first. In these two cases the saying holds.

\subsection{Limitations}
\begin{itemize}
  \item \emph{Mean flow only.} The CFD resolves the mean flow; instantaneous peaks of a meandering cloud can exceed the
  mean by an order of magnitude. This can lift the sub-threshold audible doses on the sofa (OAV 0.2--0.6) above
  threshold, but not the silent ones ($10^{-10}$). In the forward-leaning posture the mean silent cloud passes within
  25\,cm of the neighbour's face, where a filament could reach the nose.
  \item \emph{Body plume.} The seated body plume is 40--56\% too fast, which shortens the travel times of the silent
  release by about 30\% in the seated case; the arrival times in the elevator (8.8 and 17.6\,s) are likely short by a
  similar amount. Standing plumes were not validated separately.
  \item \emph{Scenes and weights.} The as-emitted shares of Figure~\ref{fig:routing} weight the simulated scenes
  equally and assume the flow path does not depend on the passage volume. They describe these scenes, not how often
  each arises.
  \item \emph{Base rate.} The share of passages smelled has never been measured; the remembered belief of
  0.67--0.92 spans the observer models that bracket it.
\end{itemize}
